% =============================================================================
% Release 2 — Sprint 1 : Mise en place de l'infrastructure
% =============================================================================

\section[Sprint 1 : Mise en place de l'infrastructure]%
{Sprint 1 : Mise en place de l'infrastructure}

Ce sprint pose les fondations cloud de \textit{TicketManager} sur Azure Stack Hub.
L'objectif est de disposer d'un réseau segmenté, de machines virtuelles d'administration
et de base de données, ainsi que d'un cluster Kubernetes opérationnel, le tout
provisionné de manière reproductible.


% ---------------------------------------------------------------------------
\subsection{Objectif du Sprint 1}
% ---------------------------------------------------------------------------

Le Sprint~1 succède au socle applicatif de la Release~1. Il se concentre sur
\textbf{l'approvisionnement de l'infrastructure} et la \textbf{préparation du cluster}
destiné à héberger l'application conteneurisée. Le tableau~\ref{tab:r2-s1-objectifs} détaille
les objectifs par axe.

\begin{table}[H]
\centering
\caption{Objectifs du Sprint~1 par axe}
\label{tab:r2-s1-objectifs}
\footnotesize
\renewcommand{\arraystretch}{1.1}
\setlength{\tabcolsep}{6pt}
\begin{tabular}{|p{2.8cm}|p{9.2cm}|}
\hline
\textbf{Axe} & \textbf{Objectif} \\
\hline
Infrastructure as Code &
  Décrire et provisionner le réseau, la sécurité et les VMs par Terraform
  (fournisseur \textit{azurestack}) \\
\hline
Segmentation réseau &
  Structurer le VNet \texttt{10.1.0.0/16} en sous-réseaux \texttt{mgmt}, \texttt{db}
  et \texttt{aks}, chacun protégé par un NSG \\
\hline
Machines virtuelles &
  Déployer \textbf{vm-gestion} (administration, agent DevOps) et \textbf{vm-db}
  (MongoDB legacy / migration) \\
\hline
Orchestration &
  Créer le cluster Kubernetes (1 master + 2 workers) via \textbf{AKS Engine} \\
\hline
Vérification &
  Contrôler les ressources dans le portail Azure Stack Hub et via \texttt{kubectl} \\
\hline
\end{tabular}
\end{table}

\paragraph{Résultat attendu.}
À l'issue du sprint :
\begin{itemize}
  \item le réseau virtuel \texttt{ticketmanager-iac-vnet} et les trois sous-réseaux sont
        créés ;
  \item \textbf{vm-gestion} (\texttt{10.1.1.4}) est accessible en SSH et héberge l'agent
        Azure DevOps, Docker, Ansible et AKS Engine ;
  \item \textbf{vm-db} (\texttt{10.1.2.4}) est joignable en IP privée depuis \texttt{mgmt}
        et \texttt{aks} ;
  \item le cluster Kubernetes affiche trois nœuds \texttt{Ready}
        (\texttt{kubectl get nodes}).
\end{itemize}

\paragraph{Périmètre.}
Le tableau~\ref{tab:r2-s1-perimetre} précise le contenu du sprint par rapport aux livraisons
ultérieures.

\begin{table}[H]
\centering
\caption{Périmètre du Sprint~1}
\label{tab:r2-s1-perimetre}
\footnotesize
\renewcommand{\arraystretch}{1.1}
\setlength{\tabcolsep}{6pt}
\begin{tabular}{|p{5.5cm}|c|c|}
\hline
\textbf{Élément} & \textbf{Sprint~1} & \textbf{Sprint ultérieur} \\
\hline
Socle applicatif (Release~1) & Livré & — \\
\hline
Terraform — VNet, NSG, NIC, VMs & Oui & — \\
\hline
AKS Engine — cluster K8s & Oui & — \\
\hline
Pipeline Azure DevOps & — & Sprint~2 \\
\hline
Déploiement pods applicatifs & — & Sprint~3 \\
\hline
Prometheus / Grafana & — & Sprint~3 \\
\hline
\end{tabular}
\end{table}


% ---------------------------------------------------------------------------
\subsection{Backlog du Sprint 1}
% ---------------------------------------------------------------------------

Le backlog du Sprint~1 est organisé en deux epics : infrastructure Terraform et cluster
Kubernetes. Les tableaux~\ref{tab:r2-s1-backlog-infra} et~\ref{tab:r2-s1-backlog-k8s} listent
les user stories planifiées.

\subsubsection{Epic — Infrastructure cloud (Terraform)}

\begin{table}[H]
\centering
\caption{Backlog Sprint~1 — Epic infrastructure Terraform}
\label{tab:r2-s1-backlog-infra}
\footnotesize
\renewcommand{\arraystretch}{0.9}
\setlength{\tabcolsep}{4pt}
\begin{tabular}{|c|p{7.8cm}|c|c|}
\hline
\textbf{ID} & \textbf{User story} & \textbf{Priorité} & \textbf{Est. (j)} \\
\hline
R2-S1-I01 &
  En tant que développeur, je provisionne l'infrastructure par Terraform sur Azure Stack Hub &
  H & 3 \\
\hline
R2-S1-I02 &
  En tant que développeur, je segmente le réseau en trois sous-réseaux dédiés (mgmt, db, aks) &
  H & 2 \\
\hline
R2-S1-I03 &
  En tant que développeur, je configure les groupes de sécurité réseau par sous-réseau &
  H & 2 \\
\hline
R2-S1-I04 &
  En tant que développeur, je dispose d'une VM d'administration (vm-gestion) pour orchestrer les déploiements &
  H & 2 \\
\hline
R2-S1-I05 &
  En tant que développeur, je déploie une VM base de données (vm-db) pour MongoDB Legacy &
  M & 1 \\
\hline
R2-S1-I06 &
  En tant que développeur, j'associe chaque VM à une interface réseau sécurisée &
  H & 1 \\
\hline
\end{tabular}
\end{table}

\subsubsection{Epic — Cluster Kubernetes (AKS Engine)}

\begin{table}[H]
\centering
\caption{Backlog Sprint~1 — Epic cluster Kubernetes}
\label{tab:r2-s1-backlog-k8s}
\footnotesize
\renewcommand{\arraystretch}{0.9}
\setlength{\tabcolsep}{4pt}
\begin{tabular}{|c|p{7.8cm}|c|c|}
\hline
\textbf{ID} & \textbf{User story} & \textbf{Priorité} & \textbf{Est. (j)} \\
\hline
R2-S1-K01 &
  En tant que développeur, je prépare le fichier de configuration AKS Engine (\texttt{kubernetes.json}) &
  H & 2 \\
\hline
R2-S1-K02 &
  En tant que développeur, je déploie le cluster Kubernetes (1 master + 2 workers) &
  H & 3 \\
\hline
R2-S1-K03 &
  En tant que développeur, je vérifie l'état des nœuds du cluster (\texttt{Ready}) &
  H & 1 \\
\hline
R2-S1-K04 &
  En tant que développeur, je configure l'accès \texttt{kubeconfig} depuis vm-gestion &
  H & 1 \\
\hline
\end{tabular}
\end{table}


% ---------------------------------------------------------------------------
\subsection{Conception de l'infrastructure}
% ---------------------------------------------------------------------------

\subsubsection{Plan d'adressage réseau}

Le tableau~\ref{tab:r2-s1-adressage} résume le plan d'adressage retenu pour l'infrastructure
TicketManager. La figure~\ref{fig:r2-s1-plan-adressage} en propose une représentation
schématique (sous-réseaux, adresses IP et séparation Terraform / AKS Engine).

\begin{table}[H]
\centering
\caption{Plan d'adressage réseau TicketManager}
\label{tab:r2-s1-adressage}
\footnotesize
\renewcommand{\arraystretch}{1.1}
\setlength{\tabcolsep}{6pt}
\begin{tabular}{|l|l|}
\hline
\textbf{Élément} & \textbf{Valeur} \\
\hline
Réseau virtuel (VNet) & \texttt{10.1.0.0/16} \\
\hline
Subnet administration (\texttt{mgmt}) & \texttt{10.1.1.0/24} \\
\hline
Subnet base de données (\texttt{db}) & \texttt{10.1.2.0/24} \\
\hline
Subnet Kubernetes (\texttt{aks}) & \texttt{10.1.3.0/24} \\
\hline
IP vm-gestion (privée) & \texttt{10.1.1.4} \\
\hline
IP vm-db (privée) & \texttt{10.1.2.4} \\
\hline
Workers K8s (exemple) & \texttt{10.1.3.4}, \texttt{10.1.3.5} \\
\hline
\end{tabular}
\end{table}

\begin{figure}[htbp]
    \centering
    \includegraphics[width=0.95\textwidth]{plan-adressage-reseau-sprint1.png}
    \caption{Plan d'adressage réseau TicketManager — vue schématique (draw.io)}
    \label{fig:r2-s1-plan-adressage}
\end{figure}

\subsubsection{Architecture globale}

L'architecture repose sur une organisation \textbf{hybride} : couche réseau et compute
provisionnée par Terraform, complétée par un cluster Kubernetes déployé via AKS Engine
sur le sous-réseau \texttt{aks}. La figure~\ref{fig:r2-s1-arch-infra} en donne une vue
d'ensemble.

\begin{figure}[htbp]
    \centering
    \includegraphics[width=0.95\textwidth]{architecture-infrastructure-sprint1.png}
    \caption{Architecture globale de l'infrastructure TicketManager (Sprint~1)}
    \label{fig:r2-s1-arch-infra}
\end{figure}

\paragraph{Composants provisionnés par Terraform.}
\begin{itemize}
  \item \textbf{Réseau virtuel et sous-réseaux} : connectivité interne entre VMs et cluster ;
  \item \textbf{Groupes de sécurité réseau (NSG)} : filtrage SSH, MongoDB, trafic interne ;
  \item \textbf{Interfaces réseau (NIC)} : rattachement des VMs au VNet ;
  \item \textbf{Machines virtuelles} : vm-gestion et vm-db
        (fichiers \texttt{vm-management.tf}, \texttt{vm-db.tf}).
\end{itemize}

\paragraph{Cluster Kubernetes (hors Terraform).}
Terraform prépare le sous-réseau \texttt{aks} et les règles NSG associées ; le déploiement
effectif du plan de contrôle et des workers est réalisé par \textbf{AKS Engine}, outil
open-source adapté à Azure Stack Hub.


% ---------------------------------------------------------------------------
\subsection{Réalisation}
% ---------------------------------------------------------------------------

\subsubsection{Infrastructure as Code avec Terraform}

Terraform adopte une approche \textbf{déclarative} : l'état cible des ressources est décrit
dans des fichiers de configuration ; le moteur calcule et applique les écarts. Le cycle
repose sur quatre commandes : \texttt{terraform init}, \texttt{validate}, \texttt{plan} et
\texttt{apply}.

La figure~\ref{fig:r2-s1-tf-init} présente le résultat de l'initialisation de l'environnement
Terraform. La figure~\ref{fig:r2-s1-tf-validate} confirme la validité syntaxique de la
configuration. La figure~\ref{fig:r2-s1-tf-plan} illustre le plan d'exécution généré avant
application. Enfin, la figure~\ref{fig:r2-s1-tf-apply} montre les sorties du déploiement,
notamment les adresses IP attribuées aux machines virtuelles.

\begin{figure}[htbp]
    \centering
    \includegraphics[width=0.85\textwidth]{s1-terraform-init.png}
    \caption{Initialisation Terraform (\texttt{terraform init})}
    \label{fig:r2-s1-tf-init}
\end{figure}

\begin{figure}[htbp]
    \centering
    \includegraphics[width=0.85\textwidth]{s1-terraform-validate.png}
    \caption{Validation Terraform (\texttt{terraform validate})}
    \label{fig:r2-s1-tf-validate}
\end{figure}

\begin{figure}[htbp]
    \centering
    \includegraphics[width=0.85\textwidth]{s1-terraform-plan.png}
    \caption{Plan d'exécution Terraform (\texttt{terraform plan})}
    \label{fig:r2-s1-tf-plan}
\end{figure}

\begin{figure}[htbp]
    \centering
    \includegraphics[width=0.85\textwidth]{s1-terraform-apply.png}
    \caption{Application Terraform — sorties (IPs vm-gestion, vm-db)}
    \label{fig:r2-s1-tf-apply}
\end{figure}

\paragraph{Vérification portail Azure Stack Hub.}
Les figures~\ref{fig:r2-s1-vnet},~\ref{fig:r2-s1-nsg} et~\ref{fig:r2-s1-vms} confirment la
création des ressources dans le groupe \texttt{AmalBH-RG}. La figure~\ref{fig:r2-s1-vnet}
présente le réseau virtuel et les sous-réseaux. La figure~\ref{fig:r2-s1-nsg} détaille les
groupes de sécurité réseau. La figure~\ref{fig:r2-s1-vms} liste les machines virtuelles
déployées.

\begin{figure}[htbp]
    \centering
    \includegraphics[width=0.85\textwidth]{s1-vnet-subnets.png}
    \caption{Réseau virtuel et sous-réseaux dans le portail}
    \label{fig:r2-s1-vnet}
\end{figure}

\begin{figure}[htbp]
    \centering
    \includegraphics[width=0.85\textwidth]{s1-nsg.png}
    \caption{Groupes de sécurité réseau}
    \label{fig:r2-s1-nsg}
\end{figure}

\begin{figure}[htbp]
    \centering
    \includegraphics[width=0.85\textwidth]{s1-vms-portal.png}
    \caption{Machines virtuelles déployées}
    \label{fig:r2-s1-vms}
\end{figure}


\subsubsection{Déploiement du cluster Kubernetes avec AKS Engine}

\paragraph{Configuration du cluster.}
Le fichier \texttt{kubernetes.json} (modèle API AKS Engine) spécifie la version de
Kubernetes, le plugin réseau (\texttt{kubenet}), le runtime (\texttt{containerd}), la
topologie des nœuds et l'authentification SSH. La figure~\ref{fig:r2-s1-aks-json} présente
un extrait de ce fichier de configuration.

\begin{figure}[htbp]
    \centering
    \includegraphics[width=0.85\textwidth]{s1-aks-engine-json.png}
    \caption{Extrait du fichier de configuration AKS Engine}
    \label{fig:r2-s1-aks-json}
\end{figure}

\paragraph{Exécution du déploiement.}
Le cluster est créé depuis vm-gestion par la commande \texttt{aks-engine deploy}. La
figure~\ref{fig:r2-s1-aks-deploy} illustre l'exécution réussie de cette commande.

\begin{figure}[htbp]
    \centering
    \includegraphics[width=0.85\textwidth]{s1-aks-engine-deploy.png}
    \caption{Déploiement du cluster avec AKS Engine}
    \label{fig:r2-s1-aks-deploy}
\end{figure}

\paragraph{Vérification.}
La figure~\ref{fig:r2-s1-kubectl-nodes} présente la sortie de \texttt{kubectl get nodes} :
les trois nœuds du cluster sont en état \texttt{Ready}.

\begin{figure}[htbp]
    \centering
    \includegraphics[width=0.85\textwidth]{s1-kubectl-nodes.png}
    \caption{Vérification des nœuds Kubernetes (\texttt{kubectl get nodes})}
    \label{fig:r2-s1-kubectl-nodes}
\end{figure}


% ---------------------------------------------------------------------------
\subsection{Bilan du Sprint 1}
% ---------------------------------------------------------------------------

Le Sprint~1 fournit l'infrastructure réseau et compute nécessaire aux sprints suivants :
vm-gestion comme point d'ancrage DevOps, sous-réseau \texttt{aks} prêt pour les pods, et
cluster Kubernetes opérationnel. Le pipeline CI/CD (Sprint~2) et le déploiement applicatif
(Sprint~3) s'appuient sur ce socle.
